\documentclass{article}
\usepackage[UTF8]{ctex}
\usepackage{amsmath,amssymb,amsfonts}
\usepackage{enumitem}
\usepackage{geometry}
\geometry{a4paper, margin=1in}
% === 错题本排版优化（不参与题目标识） ===
\setlength{\parindent}{0pt}
\setlength{\parskip}{0.35em plus 0.10em minus 0.08em}
\setlength{\emergencystretch}{3em}
\setlist{itemsep=0.25em, topsep=0.40em, parsep=0pt}
\allowdisplaybreaks[2]
% === 错题本排版优化结束 ===
\begin{document}
\section*{华南理工大学2025年数学分析真题}
\section*{156．华南理工大学2025年数学分析试卷}

\subsection{题目1}
1．（12 分）已知数列 \(x_{0}=a>0, x_{n}=\arctan x_{n-1}(n=1,2, \cdots)\) ．证明：

\begin{itemize}
  \item[（1）] \(\lim _{n \rightarrow \infty} x_{n}=0\) ．
  \item[（2）] \(\lim _{n \rightarrow \infty} \sqrt{\frac{2 n}{3}} x_{n}=1\) ．
\end{itemize}

\subsection{题目2}
1．（12 分）证明函数 \(f(x)=\left\{\begin{array}{ll}\sqrt{x} \ln x, & x>0 ; \\ 0, & x=0 .\end{array}\right.\) 在 \([0,+\infty)\) 上一致连续．\\

\subsection{题目3}
3．（12 分）证明不等式 \(\frac{1}{\sin ^{2} x} \leq \frac{1}{x^{2}}+1-\frac{4}{\pi^{2}}\left(0<x \leq \frac{\pi}{2}\right)\) ．\\

\subsection{题目4}
11．（14 分）设 \(f(x, y)=\left\{\begin{array}{ll}g(x, y) \sin \frac{1}{\sqrt{x^{2}+y^{2}}}, & (x, y) \neq(0,0) ; \\ 0, & (x, y)=(0,0) .\end{array}\right.\) 证明：

\begin{itemize}
  \item[（1）] 若 \(g(0,0)=0, g(x, y)\) 在 \((0,0)\) 处可微且 \(\mathrm{d} g(0,0)=0\) ．则 \(f(x, y)\) 在 \((0,0)\) 处可微，且 \(\mathrm{d} f(0,0)=0\) ．
  \item[（2）] 若 \(g(x, y)\) 在 \((0,0)\) 处有偏导数且 \(f(x, y)\) 在 \((0,0)\) 处可微，则 \(\mathrm{d} f(0,0)=0\) ．
\end{itemize}

\subsection{题目5}
5．（12 分）设 \(f(x)\) 在 \([0,1]\) 上连续，在 \((0,1)\) 内可导，\(f(0)=1\) ．证明：

\begin{itemize}
  \item[（1）] 存在 \(\xi \in(0,1)\) ，使得 \(f(1)=(1+\xi) f^{\prime}(\xi) \ln 2+1\) ．
  \item[（2）] \(\lim _{n \rightarrow \infty} n(\sqrt[n]{1+x}-1)=\ln (1+x), x \in(0,1)\) ．
\end{itemize}

\subsection{题目6}
6．（12 分）设 \(f(x)\) 在 \([0,1]\) 上可积，证明：存在折线函数列 \(\left\{f_{n}(x)\right\}_{n=1}^{\infty}\) ，使得

\[
\int_{0}^{1} f(x) \mathrm{d} x=\lim _{n \rightarrow \infty} \int_{0}^{1} f_{n}(x) \mathrm{d} x
\]

\subsection{题目7}
7．（12 分）求幂级数 \(\sum_{n=1}^{\infty} \frac{x^{4 n-1}}{4 n+1}\) 的收敛域与和函数 \(S(x)\) ．\\

\subsection{题目8}
8 ．（12 分）计算曲面积分 \(F(t)=\iint_{x+y+z=t} f(x, y, z) \mathrm{d} S\) ，其中

\[
f(x, y, z)= \begin{cases}1-x^{2}-y^{2}-z^{2}, & x^{2}+y^{2}+z^{2} \leq 1 \\ 0, & x^{2}+y^{2}+z^{2}>1\end{cases}
\]

\subsection{题目9}
9．（12 分）设 \(\int_{0}^{+\infty} f(x) \mathrm{d} x\) 收敛，且 \(\lim _{x \rightarrow+\infty} f(x)\) 有意义（有限数或 \(+\infty\) 或 \(-\infty\) ）．证明： \(\lim _{x \rightarrow+\infty} f(x)=0\) ．\\

\subsection{题目10}
10．（12 分）设 \(f(x)\) 在 \([0,+\infty)\) 上连续且 \(\lim _{x \rightarrow+\infty} f(x)=k\) ，试证：

\[
\int_{0}^{+\infty} \frac{f(a x)-f(b x)}{x} \mathrm{~d} x=[f(0)-k] \ln \frac{b}{a}(a>0, b>0) .
\]

\subsection{题目11}
11．（13 分）设 \(u=x+y, v=\frac{1}{x}+\frac{1}{y}\) ，试用新变量 \(u, v\) 变换等式

\[
x^{2} \frac{\partial^{2} z}{\partial x^{2}}-\left(x^{2}+y^{2}\right) \frac{\partial^{2} z}{\partial x \partial y}+y^{2} \frac{\partial^{2} z}{\partial y^{2}}=0 .
\]

（假设所有出现的二阶偏导数都连续）\\

\subsection{题目12}
12．（15 分）设函数列 \(f_{n}(x)=n x e^{-n^{2} x^{2}}, x \in[0,1], n=1,2, \cdots\) ．证明：

\begin{itemize}
  \item[（1）] \(\left\{f_{n}(x)\right\}_{n=1}^{\infty}\) 在 \([0,1]\) 上收敛于 \(f(x) \equiv 0\) ．
  \item[（2）] \(\left\{f_{n}(x)\right\}_{n=1}^{\infty}\) 在 \([0,1]\) 上是否一致收敛于 \(f(x) \equiv 0\) ？判断并给出理由．
  \item[（3）] \(\left\{f_{n}(x)\right\}_{n=1}^{\infty}\) 在 \([0,1]\) 上积分平均收敛于 \(f(x) \equiv 0\) ，即 \(\lim _{n \rightarrow \infty} \int_{0}^{1}\left|f_{n}(x)-f(x)\right| \mathrm{d} x=0\) ．
\end{itemize}
\end{document}
